\chapter{Il tessuto connettivo}

% ---------------------------------------------------------------------------
\section{Classificazione}

\textbf{Tessuto connettivo}: si distingue per l'abbondante \textbf{matrice extracellulare} che
separa le cellule. Tre grandi categorie:
\begin{itemize}
  \item \textbf{propriamente detto}:
    \begin{itemize}
      \item \textit{lasso}: fibre in rete aperta (areolare, adiposo, reticolare);
      \item \textit{denso}: fibre impacchettate (denso regolare, irregolare, elastico);
    \end{itemize}
  \item \textbf{liquido}:
    \begin{itemize}
      \item \textit{sangue}: sistema circolatorio;
      \item \textit{linfa}: sistema linfatico;
    \end{itemize}
  \item \textbf{di sostegno}:
    \begin{itemize}
      \item \textit{cartilagine}: matrice solida ed elastica (ialina, elastica, fibrosa);
      \item \textit{osso}: matrice solida a struttura cristallina.
    \end{itemize}
\end{itemize}

% ---------------------------------------------------------------------------
\section{Il sangue}

\subsection{Generalità e composizione}
\textbf{Sangue}: tessuto connettivo liquido racchiuso in un sistema di canali comunicanti (vasi
arteriosi e venosi). Fornisce nutrienti (vitamine, elettroliti) e $O_2$ alle cellule e trasporta i
prodotti di scarto agli organi emuntori.

Per centrifugazione $\rightarrow$ tre frazioni:
\begin{itemize}
  \item \textbf{plasma} (55\%): parte liquida (acqua, proteine, lipidi, glucosio, amminoacidi, ioni);
  \item \textbf{globuli bianchi e piastrine} ($\sim$1\%): strato intermedio (\textit{buffy coat});
  \item \textbf{globuli rossi} (45\%): frazione più densa (\textit{ematocrito}).
\end{itemize}
Volume: $\sim$7--8\% del peso corporeo (4,5--6 l). Le \textbf{vene} fungono da tampone: tollerano
una perdita acuta fino al 20\% ($\sim$1 l). Il plasma è in equilibrio osmotico con i liquidi
extracellulari.

\subsection{Funzioni del sangue}
\begin{itemize}
  \item \textbf{trasporto}: nutrienti, ormoni, gas disciolti;
  \item \textbf{riserva}: nutrienti e fattori della coagulazione;
  \item \textbf{termoregolazione}: redistribuzione del calore;
  \item \textbf{difesa}: contro patogeni e sostanze estranee;
  \item \textbf{antiemorragica}: coagulazione e riparazione delle lesioni vascolari;
  \item \textbf{omeostasi}: pH e composizione elettrolitica dei liquidi interstiziali.
\end{itemize}

% ---------------------------------------------------------------------------
\section{Il plasma e le proteine plasmatiche}

\subsection{Plasma e siero}
Differenza = fattori della coagulazione: quando il fibrinogeno precipita in fibrina, il plasma
diventa siero.
\begin{itemize}
  \item \textbf{siero} = albumina + globuline;
  \item \textbf{plasma} = albumina + globuline + fibrinogeno + fattori della coagulazione.
\end{itemize}

\subsection{Proteine plasmatiche}
Catene polipeptidiche, talvolta coniugate (lipoproteine, glicoproteine). Funzioni: difesa
immunitaria umorale, pressione oncotica, trasporto di sostanze insolubili in acqua, regolazione
del pH, coagulazione.

\textbf{Sintesi}: epatociti (albumina, globuline), plasmacellule (immunoglobuline), intestino
(lipoproteine). Metabolismo rapido: $\sim$9--15\% rinnovato ogni giorno.

\subsection{Frazioni elettroforetiche e valori di riferimento}
Concentrazione totale: 6--8 g/dl. Zone:
\begin{itemize}
  \item \textbf{$\alpha_1$}: antitripsina, glicoproteina acida;
  \item \textbf{$\alpha_2$}: aptoglobina, macroglobulina, antiplasmina;
  \item \textbf{$\beta$}: transferrina, LDL, complemento C3;
  \item \textbf{$\gamma$}: immunoglobuline (IgA, IgD, IgE, IgG, IgM).
\end{itemize}
Le \textbf{bande monoclonali} (picco singolo anomalo, di norma in zona $\gamma$) indicano patologie
come il mieloma multiplo.

\begin{table}[htbp]
  \centering \small \renewcommand{\arraystretch}{1.3}
  \begin{tabularx}{\textwidth}{|l|c|X|X|}
    \hline
    \textbf{Proteina} & \textbf{g/dl} & \textbf{Aumenta in} & \textbf{Diminuisce in} \\
    \hline
    Albumina & 3,6--4,9
      & disidratazione, vomito, diarrea, sudorazione eccessiva
      & malnutrizione, malattie renali/epatiche, infiammazioni, ipertiroidismo, gravidanza \\
    \hline
    $\alpha_1$-globuline & 0,2--0,4
      & infiammazioni croniche/infettive, infarto, gravidanza
      & malattie epatiche gravi, enfisema congenito, malattie renali \\
    \hline
    $\alpha_2$-globuline & 0,4--0,8
      & malattie renali e infiammatorie, infezioni, diabete, s.\ di Down
      & malattie epatiche gravi, diabete, emolisi, artrite reumatoide \\
    \hline
    $\beta$-globuline & 0,6--1,0
      & anemia sideropenica, mieloma, ipercolesterolemia, gravidanza
      & malnutrizione, cirrosi \\
    \hline
    $\gamma$-globuline & 0,9--1,4
      & gammopatie (MGUS), mieloma multiplo, epatopatie croniche, infezioni
      & malattie ereditarie del sistema immunitario \\
    \hline
  \end{tabularx}
  \caption{Frazioni proteiche plasmatiche: valori normali e cause di variazione.}
  \label{tab:proteine-plasmatiche}
\end{table}

\subsection{Albumina}
\begin{itemize}
  \item proteina plasmatica più abbondante;
  \item trasporto a bassa specificità e amplissima capacità: farmaci, bilirubina, acidi grassi,
    metalli, ormoni;
  \item dosaggio = indice della \textbf{funzionalità epatica}; diminuisce nei processi infiammatori;
  \item \textbf{analbuminemia} (assenza/grave riduzione per deficit di sintesi): senza
    sintomatologia marcata, grazie al compenso delle altre frazioni.
\end{itemize}

% ---------------------------------------------------------------------------
\section{Gli eritrociti}

\subsection{Morfologia}
\textbf{Eritrociti} (globuli rossi, \textit{RBC}): cellule anucleate, prive di organuli.
\begin{itemize}
  \item forma a \textbf{lente biconcava}: $\phi$ 7,5\,$\mu$m, spessore $\sim$2\,$\mu$m (1\,$\mu$m
    al centro) $\rightarrow$ massimizza il rapporto superficie/volume $\rightarrow$ scambio gassoso;
  \item citoplasma con \textbf{emoglobina} ed enzimi solubili, tra cui l'\textbf{anidrasi carbonica}
    (forma il bicarbonato che tampona il pH);
  \item numero: 5 mln/mm\textsuperscript{3} (maschio), 4,5 mln/mm\textsuperscript{3} (femmina);
  \item vita media 120 giorni $\rightarrow$ poi la membrana espone oligosaccaridi che la rendono
    bersaglio dei macrofagi di milza, midollo osseo, fegato.
\end{itemize}
\textbf{Membrana}: 50\% proteine (per lo più intrinseche), 40\% lipidi, 10\% carboidrati. Il
citoscheletro (spettrina, actina, anchirina, proteine di banda 3, 4.1, 4.2, 4.9) forma un guscio
unito alla membrana in molti punti $\rightarrow$ flessibilità per percorrere i capillari
(\textit{impilato} con gli altri eritrociti).

\subsection{Emoglobina}
\textbf{Emoglobina} (Hb): proteina tetramerica (p.m.\ 68.000 Da), 2 catene $\alpha$ + 2 catene
$\beta$, ciascuna con un \textbf{gruppo eme} ($Fe^{2+}$). Ogni ferro lega reversibilmente una
molecola di $O_2$:
\begin{itemize}
  \item \textbf{$HbO_2$} (ossiemoglobina): sangue arterioso, rosso vivo;
  \item \textbf{Hb} (deossiemoglobina): sangue venoso, blu-violaceo.
\end{itemize}
\textbf{Hb fetale} (HbF, catene $\alpha_2\gamma_2$): affinità per l'$O_2$ maggiore dell'Hb adulta
$\rightarrow$ favorisce il passaggio di ossigeno dalla madre al feto.

\subsection{Metabolismo del ferro}
\begin{itemize}
  \item ripartizione: $\sim$2/3 su Hb $\cdot$ 1/4 in deposito (\textbf{ferritina}) $\cdot$ resto
    funzionale (mioglobina, enzimi); totale 2 g (donna), 5 g (uomo);
  \item perdite: 2 mg/die donna, 1 mg/die uomo;
  \item assorbimento: nel \textbf{duodeno}, solo come $Fe^{2+}$;
  \item \textbf{carenza} $\rightarrow$ inibisce la sintesi di Hb (perdite di sangue, scarso apporto,
    aumentato fabbisogno in gravidanza, ridotto riciclo nelle infezioni croniche);
  \item \textbf{sovraccarico} $\rightarrow$ danno a fegato, pancreas, miocardio; se la transferrina
    è satura, il Fe libero è tossico.
\end{itemize}

\subsection{Porfirie}
\textbf{Porfirie}: malattie metaboliche rare, per lo più ereditarie $\rightarrow$ deficit di un
enzima della via dell'eme $\rightarrow$ accumulo di porfirine/precursori (escrezione in urine e
feci); manifestazioni su \textbf{sistema nervoso} e \textbf{cute} (fotosensibilità).

\subsection{Anemie emolitiche ereditarie}
\textbf{Anemie falciformi} e \textbf{sferocitosi} ereditaria: eritrociti di forma anomala
$\rightarrow$ vita media ridotta $\rightarrow$ nonostante l'eritropoiesi compensatoria, trasporto
di $O_2$ insufficiente.

\textbf{Anemia falciforme}: mutazione puntiforme nel gene della catena $\beta$ (GAG $\rightarrow$
GUG) $\rightarrow$ acido glutammico sostituito da valina $\rightarrow$ globuli rossi a falce che
ostacolano il microcircolo. Più frequente in Africa centro-orientale, sud Arabia, area mediterranea
(Italia, Spagna, Grecia). Terapie: trasfusioni, splenectomia, induttori dell'Hb fetale.

% ---------------------------------------------------------------------------
\section{Gruppi sanguigni}

\textbf{Gruppi sanguigni} (sistema AB0, Landsteiner 1901): la membrana degli eritrociti espone
catene di carboidrati ereditarie (antigeni); gli anticorpi plasmatici determinano le reazioni di
incompatibilità trasfusionale.
\begin{table}[htbp]
  \centering \small \renewcommand{\arraystretch}{1.2}
  \begin{tabular}{|c|c|l|}
    \hline
    \textbf{Gruppo} & \textbf{Antigeni} & \textbf{Anticorpi} \\
    \hline
    A  & A       & anti-B \\
    B  & B       & anti-A \\
    AB & A e B   & nessuno (ricevente universale) \\
    0  & nessuno & anti-A e anti-B (donatore universale) \\
    \hline
  \end{tabular}
\end{table}

\textbf{Fattore Rh} (isolato nel 1940 dalla scimmia \textit{Macacus rhesus}):
\begin{itemize}
  \item \textbf{Rh+}: esprime l'antigene D sulla membrana;
  \item \textbf{Rh$-$}: ne è privo, può sviluppare anticorpi anti-D dopo esposizione a sangue Rh+.
\end{itemize}

% ---------------------------------------------------------------------------
\section{I leucociti}

\subsection{Generalità}
\textbf{Leucociti} (globuli bianchi): difesa contro patogeni e sostanze estranee. Prodotti dal
midollo osseo; 5000--7000/mm\textsuperscript{3}; vita media $\sim$20 giorni.
\begin{itemize}
  \item \textbf{granulociti}: granuli colorabili, nucleo plurilobato (\textit{polimorfonucleati})
    $\rightarrow$ neutrofili, eosinofili, basofili (affiancati, per funzione, dai mastociti
    tissutali);
  \item \textbf{agranulociti}: mononucleati $\rightarrow$ linfociti e monociti (nei tessuti =
    macrofagi).
\end{itemize}
Per funzione si raggruppano anche in: \textbf{fagociti} (neutrofili, monociti/macrofagi),
\textbf{granulociti} (basofili, neutrofili, eosinofili), \textbf{cellule citotossiche} (alcuni
linfociti, incluse le natural killer), \textbf{cellule che presentano l'antigene}
(monociti/macrofagi, linfociti B/plasmacellule, cellule dendritiche).

\subsection{Granulociti neutrofili}
\begin{itemize}
  \item \textbf{50--70\%} dei leucociti; $\phi$ 12--14\,$\mu$m; nucleo plurilobato
    (\textit{polimorfonucleati});
  \item nelle donne: \textbf{corpo di Barr} (cromatina del cromosoma X inattivo);
  \item granuli con enzimi lisosomiali e sostanze battericide;
  \item primi sul sito di lesione, spiccata fagocitosi (pus), vita breve ($\le$12 h);
  \item \textbf{migrazione}: adesione all'endotelio via integrine + ICAM-1 (indotto da IL-1 e TNF)
    $\rightarrow$ diapedesi $\rightarrow$ fagocitano i batteri e rilasciano \textbf{leucotrieni}.
\end{itemize}

\subsection{Granulociti eosinofili}
\begin{itemize}
  \item $<$4\% dei leucociti;
  \item \textbf{granuli specifici}: parte interna densa (agenti antiparassitari + neurotossina),
    parte esterna meno densa;
  \item \textbf{granuli azzurrofili}: lisosomi $\rightarrow$ idrolizzano complessi antigene-anticorpo
    e parassiti fagocitati;
  \item ruolo nelle risposte anti-parassitarie e allergiche.
\end{itemize}

\subsection{Granulociti basofili}
\begin{itemize}
  \item $<$1\% dei leucociti; nucleo a S mascherato dai granuli;
  \item recettori di membrana per le IgE $\rightarrow$ \textbf{reazioni allergiche} e
    ipersensibilità immediata;
  \item vita: da 2 h a pochi giorni.
\end{itemize}

\subsection{Linfociti}
\begin{itemize}
  \item \textbf{20--25\%} dei leucociti; piccoli (8--10\,$\mu$m), nucleo grande eccentrico,
    citoplasma scarso;
  \item cellule longeve, non terminali $\rightarrow$ diventano \textbf{linfoblasti} dopo l'incontro
    con l'antigene;
  \item \textbf{linfociti T} e \textbf{B} (per marcatori di superficie) = immunità specifica;
  \item al ME: poco citoplasma, pochi mitocondri, molti ribosomi liberi, alcuni granuli azzurrofili.
\end{itemize}

\subsection{Monociti}
\begin{itemize}
  \item \textbf{1--6\%} dei leucociti; agranulociti con nucleo rotondo/reniforme;
  \item nei tessuti si differenziano in \textbf{macrofagi}: elevata fagocitosi e presentazione
    dell'antigene.
\end{itemize}

\subsection{Cellule dendritiche}
Dette anche \textbf{cellule di Langerhans} o \textbf{cellule velate}: agranulociti specializzati
nel riconoscimento dei patogeni e nell'attivazione delle altre cellule immunitarie tramite
presentazione dell'antigene.

% ---------------------------------------------------------------------------
\section{Piastrine ed emostasi}

\textbf{Piastrine} (trombociti): frammenti citoplasmatici anucleati, prodotti nel midollo osseo per
frammentazione dei \textbf{megacariociti} (cellule giganti poliploidi) $\rightarrow$ dai margini si
staccano decine di migliaia di piastrine a disco irregolare.
\begin{itemize}
  \item $\phi$ 2--4\,$\mu$m; 200.000--400.000/mm\textsuperscript{3}; vita media 8--10 giorni.
\end{itemize}

\textbf{Emostasi}: processo sequenziale che impedisce la perdita di sangue dai vasi lesi,
mantenendone integrità e fluidità. La prima fase (\textit{emostasi primaria}) = aggregazione
piastrinica nel sito di lesione. \textbf{Trombocitosi} = aumento; \textbf{trombocitopenia} =
riduzione del numero di piastrine.
